\chapter{Complexity — Universal}
%%% generated by the blueprint pipeline from TCSlib.Complexity.TuringMachine.Universal
%%%%%%%%%%%%%%%%%%%%%%%%%%%%%%%%%%%%%%%%%%%%%%%%%%%%%%%%%%%%%%%%%%%%%%%%%%%%%%%
\section{Overview}

This module provides the universal Turing machine (Arora--Barak, Theorem~1.9
and \S 1.4.1): a single machine simulating any coded machine with
multiplicative overhead depending only on the code, the quadratic corollary
for total functions of arbitrary binary machines, and a time-bounded
(clocked) universal machine that reports success or timeout against a supplied
deadline.  Most declarations build the clocked simulator: a prefix parser that
separates the clock from the code, a table interpreter with exact step counts
that stops just before each simulated transition, a fixed-width binary
countdown, and an output-buffering wrapper that withholds emissions until the
simulated machine is known to halt in budget.  Throughout, bit strings are
written over $\{0,1\}$ ($1$ for true, $0$ for false); the \emph{pair encoding}
$\langle \alpha, x\rangle$ of two bit strings lists each bit of $\alpha$
twice, then the unequal delimiter pair $01$, then $x$ verbatim; and the
\emph{buffer tape} of a string $w$ is the two-way infinite tape holding $w$ in
cells $0,\dots,|w|-1$ and blanks elsewhere.

%%%%%%%%%%%%%%%%%%%%%%%%%%%%%%%%%%%%%%%%%%%%%%%%%%%%%%%%%%%%%%%%%%%%%%%%%%%%%%%
\section{Declarations}

\begin{theorem}[The universal machine]
\lean{Turing.universal}
\label{Turing.universal}
\leanok
\uses{Turing.CodeTM.toFinTM, Turing.EffectiveMachineCode, Turing.FinTM,
  Turing.FinTM.ComputesInTime, Turing.FinTM.rightCfg,
  Turing.MultiTapeTM.runFrom, Turing.MultiTapeTM.runFrom_of_halt,
  Turing.MultiTapeTM.step_of_halt, Turing.pairEncode,
  Turing.universalBlockBound, Turing.universalCanonTM,
  Turing.universalCaptureTM, Turing.universalCapture_interpreter_run,
  Turing.universalInterpreter, Turing.universalRelation,
  Turing.universalRelation_halt, Turing.universalRelation_output,
  Turing.universalRelation_start, Turing.universalSimulationCfg,
  Turing.universalStartupBound, Turing.universalTM,
  Turing.universal_from_blocks, Turing.universal_live_block}
\difficulty{4}
For every effective machine-coding scheme there is a single finite-tape
machine $U$ over the alphabet $\{0,1\}$ such that for every code string
$\alpha$ there is a constant $C$ (depending on $\alpha$) with the following
two properties for every input string $x$.  Forward: whenever the coded
machine denoted by $\alpha$ halts on $x$ within $t$ steps with output $w$, the
machine $U$ on the pair encoding $\langle \alpha, x\rangle$ (each bit of
$\alpha$ doubled, the delimiter $01$, then $x$ verbatim) halts with the same
output $w$ within $C \cdot (t+1)$ steps.  Converse: every output that $U$
produces on $\langle \alpha, x\rangle$ upon halting is an output the denoted
machine produces on $x$ upon halting, so divergence is preserved.
\end{theorem}

\begin{theorem}[Quadratic universal simulation of total functions]
\lean{Turing.universal_quadratic}
\label{Turing.universal_quadratic}
\leanok
\uses{Turing.CodeTM.toFinTM, Turing.EffectiveMachineCode, Turing.FinTM,
  Turing.FinTM.ComputesFunInTime, Turing.FinTM.ComputesInTime,
  Turing.FinTM.ComputesInTime.mono, Turing.FinTM.one_work_tape_binary,
  Turing.MachineCode.decode_encode, Turing.exists_codeTM, Turing.pairEncode,
  Turing.universal}
\difficulty{4}
For every effective machine-coding scheme there is a single machine $U$ over
$\{0,1\}$ with the following property: whenever a finite-tape binary machine
computes a total function $f$ on bit strings within a time bound
$T \colon \mathbb{N} \to \mathbb{N}$, there are a code string $\alpha$ and a
constant $C$ such that for every input $x$, the machine $U$ on the pair
encoding $\langle \alpha, x\rangle$ halts with output $f(x)$ within
$C \cdot (T(|x|)+1)^2$ steps.
\end{theorem}

\begin{definition}[The stopped table interpreter]
\lean{Turing.timedCutInterpreter}
\label{Turing.timedCutInterpreter}
\leanok
\uses{Turing.MultiTapeTM, Turing.UniversalControl, Turing.universalInterpreter}
The four-tape table interpreter of the universal-machine construction,
modified so that every control state holding a fully decoded transition record
halts immediately instead of applying the record; on all other control states
its transition function agrees with the original interpreter.
\end{definition}

\begin{lemma}[Work-tape reads of an evaluation configuration]
\lean{Turing.timedCut_Eval_reads}
\label{Turing.timedCut_Eval_reads}
\leanok
\uses{Turing.Cfg, Turing.Cfg.workTapeSymbols, Turing.FinTM.bufferTape,
  Turing.UniversalControl, Turing.universalEvalCfg, Turing.universalFour}
\difficulty{3}
Consider an interpreter evaluation configuration, whose four work tapes are: a
table tape holding the buffer of a bit string $T$ with cursor $p$, a state
tape $\sigma \colon \mathbb{Z} \to \{0,1,\text{blank}\}$ with cursor $s$, and
a mirrored work tape and a boundary-marker tape taken from a base
configuration.  The four symbols under its work heads are, in order, the table
symbol at $p$, the state symbol $\sigma(s)$, and the two symbols read on the
corresponding tapes of the base configuration.
\end{lemma}

\begin{lemma}[Effect of one administrative action]
\lean{Turing.timedCut_Admin_apply}
\label{Turing.timedCut_Admin_apply}
\leanok
\uses{Turing.Action.apply, Turing.Cfg, Turing.UniversalControl,
  Turing.moveInputPos_zero, Turing.universalAdmin, Turing.universalEvalCfg}
\difficulty{3}
Applying an administrative interpreter action --- one that moves the table
cursor by $d_t \in \{-1,0,1\}$, optionally writes a symbol at the state-tape
cursor, moves the state cursor by $d_s$, and passes to a control state $q'$
--- to an interpreter evaluation configuration changes exactly the two
designated cursors and, when a write is present, the state-tape cell at the
cursor; the input head, the mirrored data, and the accumulated output are
unchanged.
\end{lemma}

\begin{lemma}[Administrative step rule]
\lean{Turing.timedCut_Eval_step}
\label{Turing.timedCut_Eval_step}
\leanok
\uses{Turing.Action.apply, Turing.Cfg, Turing.Cfg.inputSymbol,
  Turing.Cfg.workTapeSymbols, Turing.FinTM.bufferTape, Turing.MultiTapeTM.step,
  Turing.UniversalControl, Turing.timedCutInterpreter,
  Turing.timedCut_Admin_apply, Turing.timedCut_Eval_reads,
  Turing.universalAdmin, Turing.universalEvalCfg, Turing.universalFour}
\difficulty{3}
If, on the symbols read in an interpreter evaluation configuration, the
transition function of the stopped table interpreter returns the
administrative action with table move $d_t$, optional state-tape write, state
move $d_s$, and next control $q'$, then one machine step from that
configuration yields exactly the evaluation configuration with control $q'$,
table cursor shifted by $d_t$, and state tape and state cursor updated
accordingly.
\end{lemma}

\begin{lemma}[Reading the first unconsumed table cell]
\lean{Turing.timedCut_table_read}
\label{Turing.timedCut_table_read}
\leanok
\uses{Turing.FinTM.bufferTape, Turing.FinTM.bufferTape_nat}
\difficulty{2}
For all bit strings $l$ and $r$ and every bit $b$, the buffer tape of the
concatenation of $l$, the single bit $b$, and $r$ reads the symbol $b$ at
position $|l|$.
\end{lemma}

\begin{lemma}[Exact-cost table rewind]
\lean{Turing.timedCut_table_rewind}
\label{Turing.timedCut_table_rewind}
\leanok
\uses{Turing.Cfg, Turing.FinTM.bufferTape, Turing.FinTM.bufferTape_nat,
  Turing.MultiTapeTM.runFrom, Turing.MultiTapeTM.runFrom_succ_eq_step,
  Turing.MultiTapeTM.runFrom_zero, Turing.UniversalControl,
  Turing.timedCutInterpreter, Turing.timedCut_Eval_step,
  Turing.universalEvalCfg, Turing.universalFour, Turing.universalInterpreter}
\difficulty{4}
Suppose the stopped table interpreter is in its table-rewind phase with the
table cursor at $j-1$, where $j$ is at most the table length.  Then exactly
$j+1$ steps bring it to the count-parsing phase with the table cursor at $0$,
leaving the state tape, the mirrored data, the input head, and the accumulated
output unchanged.
\end{lemma}

\begin{lemma}[Skipping the doubled count field]
\lean{Turing.timedCut_count_run}
\label{Turing.timedCut_count_run}
\leanok
\uses{Turing.Cfg, Turing.FinTM.bufferTape, Turing.MultiTapeTM.runFrom,
  Turing.MultiTapeTM.runFrom_add, Turing.MultiTapeTM.runFrom_succ_eq_step',
  Turing.MultiTapeTM.runFrom_zero, Turing.MultiTapeTM.step,
  Turing.UniversalControl, Turing.timedCutInterpreter,
  Turing.timedCut_Eval_step, Turing.timedCut_table_read,
  Turing.universalEvalCfg, Turing.universalFour, Turing.universalInterpreter}
\difficulty{5}
Suppose the table holds, starting at the cursor position $|l|$, the doubling
$b_1 b_1 b_2 b_2 \cdots$ of a count string of $n$ bits followed by the
delimiter pair $01$.  Then the count-parsing phase of the stopped table
interpreter traverses this field in exactly $2n+2$ steps, advancing the cursor
past the delimiter and entering the initial-state copier (when copying was
requested) or the initial-state skipper; no binary arithmetic on the count's
value is performed.
\end{lemma}

\begin{lemma}[Appending a unary state symbol]
\lean{Turing.timedCut_StateTape_append}
\label{Turing.timedCut_StateTape_append}
\leanok
\uses{Turing.FinTM.bufferTape_append, Turing.universalStateTape}
\difficulty{2}
Writing the symbol $1$ into cell $n+1$ of the unary state tape representing
$n$ (a marker $0$ in cell $0$ followed by $n$ ones) yields the unary state
tape representing $n+1$.
\end{lemma}

\begin{lemma}[Blank after the last unary state symbol]
\lean{Turing.timedCut_StateTape_end}
\label{Turing.timedCut_StateTape_end}
\leanok
\uses{Turing.FinTM.bufferTape, Turing.universalStateTape}
\difficulty{1}
The unary state tape representing $n$ (a marker $0$ in cell $0$ followed by
$n$ ones) reads a blank at cell $n+1$.
\end{lemma}

\begin{lemma}[Marker-directed state rewind]
\lean{Turing.timedCut_state_rewind}
\label{Turing.timedCut_state_rewind}
\leanok
\uses{Turing.Cfg, Turing.MultiTapeTM.runFrom,
  Turing.MultiTapeTM.runFrom_succ_eq_step, Turing.MultiTapeTM.runFrom_zero,
  Turing.UniversalControl, Turing.timedCutInterpreter,
  Turing.timedCut_Eval_step, Turing.universalAdmin, Turing.universalEvalCfg,
  Turing.universalFour}
\difficulty{4}
Suppose the state tape holds the marker symbol $0$ at cell $0$ and no cell at
a positive position holds $0$, and suppose the current control of the stopped
table interpreter moves the state head left on every non-marker read while, on
reading the marker, it moves right once and transfers to a designated
continuation control $q'$.  Then from state cursor $j \in \mathbb{N}$ the
interpreter reaches $q'$ with state cursor $1$ in exactly $j+1$ steps, without
changing any tape contents, the input head, or the output; the premise is
independent of whether the traversed cells are erased blanks or retained unary
symbols.
\end{lemma}

\begin{lemma}[Skipping the initial-state field]
\lean{Turing.timedCut_initial_skip}
\label{Turing.timedCut_initial_skip}
\leanok
\uses{Turing.Cfg, Turing.FinTM.bufferTape, Turing.MultiTapeTM.runFrom,
  Turing.MultiTapeTM.runFrom_succ_eq_step, Turing.MultiTapeTM.runFrom_zero,
  Turing.UniversalControl, Turing.timedCutInterpreter,
  Turing.timedCut_Eval_step, Turing.universalEvalCfg, Turing.universalFour,
  Turing.universalInterpreter}
\difficulty{4}
If the table holds $n$ ones followed by a zero starting at the cursor position
$|l|$, then the initial-state-skipping phase of the stopped table interpreter
traverses this unary field in exactly $n+1$ steps and enters record-group
selection with the table cursor at $|l|+n+1$.
\end{lemma}

\begin{lemma}[Copying a unary field onto the state tape]
\lean{Turing.timedCut_unary_copy}
\label{Turing.timedCut_unary_copy}
\leanok
\uses{Turing.Cfg, Turing.FinTM.bufferTape, Turing.MultiTapeTM.runFrom,
  Turing.MultiTapeTM.runFrom_succ_eq_step, Turing.MultiTapeTM.runFrom_zero,
  Turing.UniversalControl, Turing.timedCutInterpreter,
  Turing.timedCut_Eval_step, Turing.timedCut_StateTape_append,
  Turing.universalAdmin, Turing.universalEvalCfg, Turing.universalFour,
  Turing.universalStateTape}
\difficulty{4}
Suppose the current control $q$ of the stopped table interpreter appends one
unary symbol to the state tape and advances both cursors whenever the table
reads $1$, and on reading $0$ transfers to a continuation control $q'$ with a
designated final table move.  If the table holds $n$ ones followed by a zero
at cursor position $|l|$, then $n+1$ steps transform the unary state tape for
$j$ (state cursor $j+1$) into the unary state tape for $j+n$ (state cursor
$j+n+1$), with the table cursor ending at $|l|+n$ plus the final move.  This
single gadget serves both initial-state extraction and successor-state
replacement.
\end{lemma}

\begin{lemma}[Installing a marker on a blank tape]
\lean{Turing.timedCut_install_marker}
\label{Turing.timedCut_install_marker}
\leanok
\uses{Turing.FinTM.bufferTape, Turing.FinTM.bufferTape_append}
\difficulty{1}
Writing a single symbol $b$ into cell $0$ of the all-blank tape yields the
buffer tape of the one-character string $b$.
\end{lemma}

\begin{definition}[Interpreter entry configuration]
\lean{Turing.timedCut_InterpreterInitial}
\label{Turing.timedCut_InterpreterInitial}
\leanok
\uses{Turing.Cfg, Turing.FinTM.bufferTape, Turing.UniversalControl,
  Turing.universalFour}
The configuration of the stopped table interpreter at entry: control in the
start phase, the captured table held on the first work tape as a buffer with
its cursor at the right blank, three blank auxiliary tapes, empty output, and
the input head already parked at a given position of the input.
\end{definition}

\begin{definition}[Inactive base data during initialization]
\lean{Turing.timedCut_InterpreterBase}
\label{Turing.timedCut_InterpreterBase}
\leanok
\uses{Turing.Cfg, Turing.FinTM.bufferTape, Turing.UniversalControl,
  Turing.universalFour}
The inactive part of an interpreter configuration during initialization: the
input head parked at a given position, a blank mirrored work tape, the
virtual-input left-boundary marker installed at cell $0$ of its tape with its
head at $1$ (also correct for an empty input suffix), and empty output.
\end{definition}

\begin{lemma}[First interpreter step]
\lean{Turing.timedCut_Interpreter_first}
\label{Turing.timedCut_Interpreter_first}
\leanok
\uses{Turing.MultiTapeTM.step, Turing.moveInputPos_zero,
  Turing.timedCutInterpreter, Turing.timedCut_InterpreterBase,
  Turing.timedCut_InterpreterInitial, Turing.timedCut_install_marker,
  Turing.universalEvalCfg, Turing.universalStateTape}
\difficulty{3}
The first step of the stopped table interpreter from its entry configuration
installs the two permanent markers (the state-tape origin marker and the
virtual-input boundary marker) and enters the unconditional table rewind with
the table cursor on the last table cell.
\end{lemma}

\begin{lemma}[Exact interpreter initialization]
\lean{Turing.timedCut_Interpreter_initialize}
\label{Turing.timedCut_Interpreter_initialize}
\leanok
\uses{Turing.MultiTapeTM.runFrom, Turing.MultiTapeTM.runFrom_add,
  Turing.MultiTapeTM.step, Turing.pairEncode, Turing.timedCutInterpreter,
  Turing.timedCut_InterpreterBase, Turing.timedCut_InterpreterInitial,
  Turing.timedCut_Interpreter_first, Turing.timedCut_count_run,
  Turing.timedCut_state_rewind, Turing.timedCut_table_rewind,
  Turing.timedCut_unary_copy, Turing.universalEvalCfg,
  Turing.universalInterpreter, Turing.universalStateTape,
  Turing.universalStateTape_marker, Turing.universal_pair_length}
\difficulty{5}
Suppose the captured table $T$ is the pair encoding of a count string of $m$
bits with a body consisting of $n$ ones, a terminator zero, and a record
region.  Then exactly $|T| + 2m + 2n + 7$ steps of the stopped table
interpreter bring the entry configuration to the main simulation phase, with
the table cursor at $2m + 2 + n + 1$ (just past the header) and the state tape
holding the unary representation of $n$ with cursor $1$; every intermediate
configuration keeps the input head fixed and the output empty, and no
transition-record lookup is involved.
\end{lemma}

\begin{lemma}[Skipping the fixed action fields]
\lean{Turing.timedCut_skip_fixed}
\label{Turing.timedCut_skip_fixed}
\leanok
\uses{Turing.Cfg, Turing.MultiTapeTM.runFrom,
  Turing.MultiTapeTM.runFrom_succ_eq_step, Turing.MultiTapeTM.runFrom_zero,
  Turing.UniversalControl, Turing.timedCutInterpreter,
  Turing.timedCut_Eval_step, Turing.universalEvalCfg,
  Turing.universalInterpreter}
\difficulty{4}
In the fixed-field-skipping phase of the stopped table interpreter, with
current field index $i$ and $n$ further fields remaining (so $i + n = 7$),
exactly $n+1$ steps advance the table cursor by $n+1$ cells and enter the
unary-tail scanner; no tape content is inspected or modified.
\end{lemma}

\begin{lemma}[Skipping a record's unary tail]
\lean{Turing.timedCut_skip_unary}
\label{Turing.timedCut_skip_unary}
\leanok
\uses{Turing.Cfg, Turing.FinTM.bufferTape, Turing.MultiTapeTM.runFrom,
  Turing.MultiTapeTM.runFrom_succ_eq_step, Turing.MultiTapeTM.runFrom_zero,
  Turing.UniversalControl, Turing.timedCutInterpreter,
  Turing.timedCut_Eval_step, Turing.universalEvalCfg, Turing.universalFour,
  Turing.universalInterpreter, Turing.universalSkipDone}
\difficulty{4}
If the table holds $n$ ones followed by a zero at cursor position $|l|$, the
unary-skipping phase of the stopped table interpreter traverses them in
exactly $n+1$ steps, ending with the cursor at $|l|+n+1$; on the terminator it
completes the skip request when the bounded record counter is zero, and
otherwise decrements the counter and re-enters fixed-field skipping.
\end{lemma}

\begin{lemma}[Skipping one serialized record]
\lean{Turing.timedCut_skip_record}
\label{Turing.timedCut_skip_record}
\leanok
\uses{Turing.Action, Turing.Cfg, Turing.MultiTapeTM.runFrom,
  Turing.MultiTapeTM.runFrom_add, Turing.UniversalControl,
  Turing.timedCutInterpreter, Turing.timedCut_skip_fixed,
  Turing.timedCut_skip_unary, Turing.universalActionBits,
  Turing.universalEvalCfg, Turing.universalNextOnes,
  Turing.universalRecordBits, Turing.universalSkipDone,
  Turing.universal_record_shape}
\difficulty{4}
If a complete serialized transition record --- eight fixed bits followed by a
unary successor field and its terminator --- occupies the table starting at
the cursor, then the skipping phase of the stopped table interpreter traverses
it in exactly its serialized length, decrementing the bounded record counter,
or completing the skip request when the counter is zero, with the cursor moved
past the record.
\end{lemma}

\begin{lemma}[Skipping a list of records]
\lean{Turing.timedCut_skip_records}
\label{Turing.timedCut_skip_records}
\leanok
\uses{Turing.Action, Turing.Cfg, Turing.MultiTapeTM.runFrom,
  Turing.MultiTapeTM.runFrom_add, Turing.UniversalControl,
  Turing.timedCutInterpreter, Turing.timedCut_skip_record,
  Turing.universalEvalCfg, Turing.universalRecordBits,
  Turing.universalSkipDone}
\difficulty{4}
If the table holds, at the cursor, the concatenated serializations of a
nonempty list of transition records whose length is one more than the bounded
record counter, then the skipping phase of the stopped table interpreter
traverses all of them in exactly the total serialized length and finishes in
the requested continuation with the cursor past the last record; no extra
transition occurs between consecutive records.
\end{lemma}

\begin{lemma}[Skipping state groups while erasing the unary state]
\lean{Turing.timedCut_skip_groups}
\label{Turing.timedCut_skip_groups}
\leanok
\uses{Turing.Action, Turing.Cfg, Turing.MultiTapeTM.runFrom,
  Turing.MultiTapeTM.runFrom_add, Turing.MultiTapeTM.runFrom_succ_eq_step,
  Turing.MultiTapeTM.runFrom_succ_eq_step', Turing.MultiTapeTM.runFrom_zero,
  Turing.UniversalControl, Turing.timedCutInterpreter,
  Turing.timedCut_Eval_step, Turing.timedCut_skip_records,
  Turing.universalEvalCfg, Turing.universalFour, Turing.universalInterpreter,
  Turing.universalRecordBits, Turing.universalSkipDone,
  Turing.universalStateWindow, Turing.universalStateWindow_end,
  Turing.universalStateWindow_erase, Turing.universalStateWindow_read}
\difficulty{5}
Suppose the table holds, at the cursor, the concatenated records of a list of
$g$ state groups, each consisting of nine records, and the state tape holds a
partially erased unary state whose remaining ones number exactly $g$.  Then in
exactly $g$ plus the total serialized length plus one steps, the stopped table
interpreter erases one unary state symbol per group, skips all of the groups'
records, and enters the state rewind, with the table cursor past the traversed
records and all erasures recorded on the state tape.
\end{lemma}

\begin{lemma}[Reading the eight fixed action bits]
\lean{Turing.timedCut_read_fixed}
\label{Turing.timedCut_read_fixed}
\leanok
\uses{Turing.Cfg, Turing.FinTM.bufferTape, Turing.FinTM.bufferTape_nat,
  Turing.MultiTapeTM.runFrom, Turing.MultiTapeTM.runFrom_succ_eq_step,
  Turing.MultiTapeTM.runFrom_zero, Turing.UniversalControl,
  Turing.timedCutInterpreter, Turing.timedCut_Eval_step,
  Turing.universalEvalCfg, Turing.universalFour, Turing.universalInterpreter}
\difficulty{5}
Suppose the table holds a word $b$ of eight bits starting at position $|l|$,
the action-reading phase of the stopped table interpreter is at field index
$i$ with $n = 7-i$ fields left, and the accumulated eight-bit register agrees
with $b$ on all indices below $i$.  Then exactly $n+1$ steps fill the register
with all of $b$ and enter the successor-decoding phase, with the table cursor
at $|l|+8$.
\end{lemma}

\begin{definition}[Successor-preparation cost]
\lean{Turing.timedCut_NextCost}
\label{Turing.timedCut_NextCost}
\leanok
The number of interpreter steps used to decode a record's successor field and
install and rewind its unary state before the action is applied: $1$ for a
halting successor, and $2q + 4$ for a live successor of index $q$.
\end{definition}

\begin{lemma}[Preparing the successor state]
\lean{Turing.timedCut_prepare_next}
\label{Turing.timedCut_prepare_next}
\leanok
\uses{Turing.Cfg, Turing.FinTM.bufferTape, Turing.MultiTapeTM.runFrom,
  Turing.MultiTapeTM.runFrom_add, Turing.MultiTapeTM.runFrom_succ_eq_step,
  Turing.MultiTapeTM.runFrom_succ_eq_step', Turing.MultiTapeTM.runFrom_zero,
  Turing.UniversalControl, Turing.timedCutInterpreter,
  Turing.timedCut_Eval_step, Turing.timedCut_NextCost,
  Turing.timedCut_state_rewind, Turing.timedCut_unary_copy,
  Turing.universalEvalCfg, Turing.universalFour, Turing.universalInterpreter,
  Turing.universalNextOnes, Turing.universalStateTape,
  Turing.universalStateTape_marker}
\difficulty{5}
Suppose the table holds, at the cursor, the unary successor field of a record
(its number of ones encoding the optional successor state) followed by its
terminator, and the state tape is the unary tape for $0$ with cursor $1$.
Then within exactly the successor-preparation cost, the stopped table
interpreter installs the unary representation of the successor index (or of
$0$ in the halting case) on the state tape with cursor $1$ and enters the
pending-action phase carrying the eight decoded bits and the halting flag,
with the table cursor stopped on the field's terminator.
\end{lemma}

\begin{definition}[The nine actions of a state]
\lean{Turing.timedCut_Actions}
\label{Turing.timedCut_Actions}
\leanok
\uses{Turing.Action, Turing.CodeTM}
For a coded one-work-tape binary machine and a state $q$, the list of the nine
transition actions of $q$, one for each pair of an optional input read and an
optional work read, ordered with the input read (blank, $0$, $1$) as the major
key and the work read (in the same order) as the minor key.
\end{definition}

\begin{lemma}[Selecting an action by read offset]
\lean{Turing.timedCut_Actions_lookup}
\label{Turing.timedCut_Actions_lookup}
\leanok
\uses{Turing.CodeTM, Turing.timedCut_Actions, Turing.timedCut_Header,
  Turing.universalRecordIndex}
\difficulty{2}
For every coded machine, state $q$, and pair of optional input and work reads,
the nine-action list of $q$ has length $9$, and its entry at the record index
determined by the two reads (a number in $\{0,\dots,8\}$) is exactly the
machine's transition action for $q$ on those reads.
\end{lemma}

\begin{definition}[Serialization header]
\lean{Turing.timedCut_Header}
\label{Turing.timedCut_Header}
\leanok
\uses{Turing.CodeTM, Turing.pairEncode}
The header of a coded machine's serialization: the doubled binary state count
with its delimiter pair, followed by the unary initial-state field (as many
ones as the initial state's index) and its terminator zero; the transition
records are excluded.
\end{definition}

\begin{lemma}[Serialization as header plus grouped records]
\lean{Turing.timedCut_serialization_actions}
\label{Turing.timedCut_serialization_actions}
\leanok
\uses{Turing.CodeTM, Turing.CodeTM.serialize, Turing.pairEncode,
  Turing.timedCut_Actions, Turing.timedCut_Header,
  Turing.universalRecordBits, Turing.universalRecords}
\difficulty{5}
The canonical serialization of a coded machine equals its header followed by
the concatenation, over its states in increasing order, of the serialized
nine-action record lists of each state.
\end{lemma}

\begin{lemma}[Decomposing the table at a selected action]
\lean{Turing.timedCut_lookup_parts}
\label{Turing.timedCut_lookup_parts}
\leanok
\uses{Turing.Action, Turing.CodeTM, Turing.CodeTM.serialize,
  Turing.timedCut_Actions, Turing.timedCut_Actions_lookup,
  Turing.timedCut_Header, Turing.timedCut_serialization_actions,
  Turing.universalRecordBits, Turing.universalRecordIndex}
\difficulty{5}
For every coded machine, state $q$, and pair of optional input and work reads,
the serialization splits as: the header, then the record groups of the $q$
states preceding $q$ (each a list of nine records), then the records of $q$'s
group preceding the read offset selected by the two reads, then the serialized
record of the machine's transition action for $q$ on those reads, then a
remainder.
\end{lemma}

\begin{lemma}[Decoding the fixed action bits]
\lean{Turing.timedCut_ActionBits_decode}
\label{Turing.timedCut_ActionBits_decode}
\leanok
\uses{Turing.Action, Turing.universalActionBits, Turing.universalSign,
  Turing.universalWrite}
\difficulty{2}
Decoding the four fixed two-bit pairs of a transition action's eight-bit
serialization recovers, in order, the action's input-head movement, its
optional work-tape write, its work-head movement, and its optional emission.
\end{lemma}

\begin{lemma}[Selecting the source record]
\lean{Turing.timedCut_select}
\label{Turing.timedCut_select}
\leanok
\uses{Turing.Action, Turing.Cfg, Turing.MultiTapeTM.runFrom,
  Turing.MultiTapeTM.runFrom_add, Turing.UniversalControl,
  Turing.timedCutInterpreter, Turing.timedCut_skip_groups,
  Turing.timedCut_skip_records, Turing.timedCut_state_rewind,
  Turing.universalEvalCfg, Turing.universalInterpreter,
  Turing.universalRecordBits, Turing.universalSkipDone,
  Turing.universalStateTape, Turing.universalStateTape_marker,
  Turing.universalStateWindow_empty, Turing.universalStateWindow_zero}
\difficulty{5}
Suppose the table holds, at the cursor, the concatenated records of $g$
preceding state groups (nine records each) followed by the records preceding
the selected read offset, and the state tape is the unary tape for $g$ with
cursor $1$.  Then in exactly $2g$ plus the total serialized length of the
traversed records plus $3$ steps, the stopped table interpreter destructively
counts down the unary state while skipping its groups, rewinds the erased
state tape, skips the read-offset prefix (entering the action reader directly
when the offset is zero), and arrives in the action-reading phase with a
cleared register, the table cursor at the selected record, and the state tape
reset to the unary tape for $0$.
\end{lemma}

\begin{lemma}[Concatenating two run equalities]
\lean{Turing.timedCut_run_join}
\label{Turing.timedCut_run_join}
\leanok
\uses{Turing.Cfg, Turing.MultiTapeTM, Turing.MultiTapeTM.runFrom,
  Turing.MultiTapeTM.runFrom_add}
\difficulty{1}
For any multi-tape machine, if running $s$ steps from a configuration $a$
yields $b$ and running $t$ steps from $b$ yields $c$, then running $s+t$ steps
from $a$ yields $c$.
\end{lemma}

\begin{lemma}[Applying the decoded record]
\lean{Turing.timed_apply_record}
\label{Turing.timed_apply_record}
\leanok
\uses{Turing.Action, Turing.Action.apply, Turing.Cfg, Turing.Cfg.inputSymbol,
  Turing.Cfg.workTapeSymbols, Turing.CodeTM, Turing.CodeTM.serialize,
  Turing.FinTM.bufferTape, Turing.FinTM.virtualMove, Turing.MultiTapeTM.step,
  Turing.UniversalControl, Turing.timedCut_ActionBits_decode,
  Turing.universalActionBits, Turing.universalEvalCfg, Turing.universalFour,
  Turing.universalInput_move, Turing.universalInput_read,
  Turing.universalInterpreter, Turing.universalSimulationCfg,
  Turing.universalStateTape}
\difficulty{5}
Let $M$ be a coded machine, let $\mathrm{src}$ be a configuration of $M$, and
let $a$ be a transition action of $M$.  One step of the (unstopped) four-tape
table interpreter from the pending-action evaluation configuration --- built
over the interpreter checkpoint of $\mathrm{src}$, carrying the eight decoded
bits of $a$ and the halting flag of its successor, with the serialized table,
any table cursor $p$, and the unary tape of the successor index installed ---
equals exactly the interpreter checkpoint of the configuration obtained by
applying $a$ to $\mathrm{src}$, with table cursor $p$; this includes the
correct virtual-input head movement, optional write, emission, and the halting
case.
\end{lemma}

\begin{lemma}[One live simulation block]
\lean{Turing.timedCut_live_block}
\label{Turing.timedCut_live_block}
\leanok
\uses{Turing.Action.apply, Turing.Cfg, Turing.Cfg.inputSymbol,
  Turing.Cfg.workTapeSymbols, Turing.CodeTM, Turing.CodeTM.serialize,
  Turing.MultiTapeTM.runFrom, Turing.MultiTapeTM.runFrom_succ_eq_step',
  Turing.MultiTapeTM.runFrom_zero, Turing.MultiTapeTM.step,
  Turing.UniversalControl, Turing.pairEncode, Turing.timedClockPrefix,
  Turing.timedCutInterpreter, Turing.timedCut_Eval_step,
  Turing.timedCut_Header, Turing.timedCut_NextCost,
  Turing.timedCut_count_run, Turing.timedCut_initial_skip,
  Turing.timedCut_lookup_parts, Turing.timedCut_prepare_next,
  Turing.timedCut_read_fixed, Turing.timedCut_run_join,
  Turing.timedCut_select, Turing.timedCut_table_rewind,
  Turing.timed_apply_record, Turing.universalActionBits,
  Turing.universalAdmin, Turing.universalEvalCfg,
  Turing.universalInput_read, Turing.universalInterpreter,
  Turing.universalNextOnes, Turing.universalRecordBits,
  Turing.universalRecordIndex, Turing.universalSimulationCfg,
  Turing.universalStateTape, Turing.universal_pair_length,
  Turing.universal_record_shape}
\difficulty{6}
Let $M$ be a coded machine with serialization of length $L$ and state count
$N+1$, let $\mathrm{src}$ be a non-halted configuration of $M$, and let the
table cursor $p$ satisfy $p \le L$.  Then there exist a duration
$d \le 3L + 5(N+1) + 20$, a new cursor $p' \le L$, and a configuration
$\mathrm{ready}$ whose control is a pending-action state, such that the
stopped table interpreter runs from the interpreter checkpoint of
$\mathrm{src}$ to $\mathrm{ready}$ in exactly $d$ steps, and one step of the
unstopped interpreter from $\mathrm{ready}$ is the interpreter checkpoint of
the configuration obtained by one step of $M$ from $\mathrm{src}$, with table
cursor $p'$.
\end{lemma}

\begin{definition}[Physical clock prefix]
\lean{Turing.timedClockPrefix}
\label{Turing.timedClockPrefix}
\leanok
The physical input region occupied by a clock string $b$ under the nested pair
encoding: each bit of $b$ repeated four times, followed by the four delimiter
cells $0011$.
\end{definition}

\begin{lemma}[Nested pairing splits off the clock]
\lean{Turing.timed_input_layout}
\label{Turing.timed_input_layout}
\leanok
\uses{Turing.pairEncode, Turing.timedClockPrefix}
\difficulty{3}
For all bit strings $b$, $\alpha$, $x$, the nested pair encoding
$\langle \langle b, \alpha\rangle, x\rangle$ equals the physical clock prefix
of $b$ (each bit quadrupled, then the four cells $0011$) followed by the pair
encoding $\langle \alpha, x\rangle$.
\end{lemma}

\begin{lemma}[Length of the clock prefix]
\lean{Turing.timedClockPrefix_length}
\label{Turing.timedClockPrefix_length}
\leanok
\uses{Turing.timedClockPrefix}
\difficulty{3}
The physical clock prefix of a string $b$ has length $4|b| + 4$.
\end{lemma}

\begin{definition}[Prefix-parser control states]
\lean{Turing.TimedPrefixControl}
\label{Turing.TimedPrefixControl}
\leanok
The finite control of the clock--code prefix parser: four phases that walk the
four physical copies of each clock bit, an end-of-clock phase for the
separator, and two phases that walk the doubled code bits.
\end{definition}

\begin{definition}[The prefix parser]
\lean{Turing.timedPrefixTM}
\label{Turing.timedPrefixTM}
\leanok
\uses{Turing.FinTM, Turing.TimedPrefixControl}
A one-work-tape machine over $\{0,1\}$ that scans the nested pair encoding
from the left: it stores each undoubled clock bit on its work tape (one write
per four input cells), then emits each undoubled code bit as output (one
emission per two input cells), and halts on the outer separator, before
reading any of the input suffix.
\end{definition}

\begin{definition}[Prefix-parser configuration]
\lean{Turing.timedPrefixCfg}
\label{Turing.timedPrefixCfg}
\leanok
\uses{Turing.Cfg, Turing.FinTM.bufferTape, Turing.TimedPrefixControl,
  Turing.pairEncode}
The configuration of the prefix parser on the nested pair encoding
$\langle \langle b, \alpha\rangle, x\rangle$ with a given control state and
input position, the captured clock string held on its work tape as a buffer
with head at the right blank, and a given accumulated output.
\end{definition}

\begin{lemma}[Length of the nested pair encoding]
\lean{Turing.timed_input_length}
\label{Turing.timed_input_length}
\leanok
\uses{Turing.pairEncode, Turing.universal_pair_length}
\difficulty{2}
For all bit strings $b$, $\alpha$, $x$, the nested pair encoding
$\langle \langle b, \alpha\rangle, x\rangle$ has length
$4|b| + 2|\alpha| + 6 + |x|$.
\end{lemma}

\begin{lemma}[Cells of a quadrupled clock bit]
\lean{Turing.timed_clock_get}
\label{Turing.timed_clock_get}
\leanok
\uses{Turing.pairEncode, Turing.timedClockPrefix, Turing.timed_input_layout}
\difficulty{4}
For every $j < |b|$ and every $r < 4$, cell $4j + r$ of the nested pair
encoding $\langle \langle b, \alpha\rangle, x\rangle$ holds the clock bit
$b_j$.
\end{lemma}

\begin{lemma}[The doubled inner separator]
\lean{Turing.timed_clock_separator}
\label{Turing.timed_clock_separator}
\leanok
\uses{Turing.pairEncode, Turing.timedClockPrefix, Turing.timed_input_layout}
\difficulty{3}
For every $r < 4$, cell $4|b| + r$ of the nested pair encoding
$\langle \langle b, \alpha\rangle, x\rangle$ holds the $r$-th symbol of the
four-cell separator $0011$.
\end{lemma}

\begin{lemma}[One advancing parser step]
\lean{Turing.timedPrefix_advance}
\label{Turing.timedPrefix_advance}
\leanok
\uses{Turing.Action.apply, Turing.Cfg.inputSymbol,
  Turing.FinTM.inputSymbol_at, Turing.MultiTapeTM.step,
  Turing.TimedPrefixControl, Turing.moveInputPos_pos_of_ne_right,
  Turing.pairEncode, Turing.timedPrefixCfg, Turing.timedPrefixTM}
\difficulty{3}
Suppose the prefix parser's transition in control $q$, on the input symbol
found at physical position $p$ of the nested pair encoding, moves the input
head right, writes nothing to the work tape, optionally emits one bit, and
passes to a control $q'$.  Then one machine step from the corresponding parser
configuration advances the physical input position by one, appends the
optional emission to the output, keeps the stored clock unchanged, and enters
$q'$.
\end{lemma}

\begin{lemma}[The clock-writing parser step]
\lean{Turing.timedPrefix_write}
\label{Turing.timedPrefix_write}
\leanok
\uses{Turing.FinTM.bufferTape_append, Turing.MultiTapeTM.step,
  Turing.moveInputPos_pos_of_ne_right, Turing.pairEncode,
  Turing.timedPrefixCfg, Turing.timedPrefixTM}
\difficulty{3}
In the fourth phase of a clock bit $b$, one step of the prefix parser advances
the input head one cell and appends the single bit $b$ to the clock string
stored on its work tape, returning to the first clock phase with no output.
\end{lemma}

\begin{lemma}[Clock extraction]
\lean{Turing.timedPrefix_clock}
\label{Turing.timedPrefix_clock}
\leanok
\uses{Turing.MultiTapeTM.initCfg, Turing.MultiTapeTM.runFrom,
  Turing.MultiTapeTM.runFrom_succ_eq_step', Turing.pairEncode,
  Turing.timedPrefixCfg, Turing.timedPrefixTM, Turing.timedPrefix_advance,
  Turing.timedPrefix_write, Turing.timed_clock_get,
  Turing.timed_input_length}
\difficulty{4}
For every $j \le |b|$, running the prefix parser $4j$ steps from its initial
configuration on the nested pair encoding
$\langle \langle b, \alpha\rangle, x\rangle$ stores exactly the first $j$
clock bits on its work tape, leaves the output empty, and places the input
head just past the $4j$ consumed cells, still in the first clock phase.
\end{lemma}

\begin{lemma}[Consuming the clock delimiter]
\lean{Turing.timedPrefix_clock_end}
\label{Turing.timedPrefix_clock_end}
\leanok
\uses{Turing.MultiTapeTM.initCfg, Turing.MultiTapeTM.runFrom,
  Turing.MultiTapeTM.runFrom_succ_eq_step', Turing.pairEncode,
  Turing.timedPrefixCfg, Turing.timedPrefixTM, Turing.timedPrefix_advance,
  Turing.timedPrefix_clock, Turing.timed_clock_separator,
  Turing.timed_input_length}
\difficulty{4}
Running the prefix parser $4|b| + 4$ steps from its initial configuration on
the nested pair encoding $\langle \langle b, \alpha\rangle, x\rangle$ consumes
the entire quadrupled clock and its four separator cells, stores exactly $b$
on the work tape, produces no output, and enters the code-extraction phase.
\end{lemma}

\begin{lemma}[Indexing past the clock prefix]
\lean{Turing.timed_code_get}
\label{Turing.timed_code_get}
\leanok
\uses{Turing.pairEncode, Turing.timedClockPrefix_length,
  Turing.timed_input_layout}
\difficulty{2}
For every $j$, cell $4|b| + 4 + j$ of the nested pair encoding
$\langle \langle b, \alpha\rangle, x\rangle$ equals cell $j$ of the inner pair
encoding $\langle \alpha, x\rangle$.
\end{lemma}

\begin{lemma}[Cells of a doubled code bit]
\lean{Turing.timed_pair_get}
\label{Turing.timed_pair_get}
\leanok
\uses{Turing.pairEncode}
\difficulty{3}
For every $j < |\alpha|$, cells $2j$ and $2j+1$ of the pair encoding
$\langle \alpha, x\rangle$ both hold the code bit $\alpha_j$.
\end{lemma}

\begin{lemma}[The aligned separator]
\lean{Turing.timed_pair_separator}
\label{Turing.timed_pair_separator}
\leanok
\uses{Turing.pairEncode}
\difficulty{3}
Cells $2|\alpha|$ and $2|\alpha| + 1$ of the pair encoding
$\langle \alpha, x\rangle$ hold $0$ and $1$ respectively.
\end{lemma}

\begin{lemma}[Code extraction]
\lean{Turing.timedPrefix_code}
\label{Turing.timedPrefix_code}
\leanok
\uses{Turing.MultiTapeTM.initCfg, Turing.MultiTapeTM.runFrom,
  Turing.MultiTapeTM.runFrom_succ_eq_step', Turing.pairEncode,
  Turing.timedPrefixCfg, Turing.timedPrefixTM, Turing.timedPrefix_advance,
  Turing.timedPrefix_clock_end, Turing.timed_code_get,
  Turing.timed_input_length, Turing.timed_pair_get}
\difficulty{4}
For every $j \le |\alpha|$, running the prefix parser $4|b| + 4 + 2j$ steps
from its initial configuration on the nested pair encoding
$\langle \langle b, \alpha\rangle, x\rangle$ keeps the stored clock $b$
intact, emits exactly the first $j$ code bits as output, and remains in the
code-extraction phase with the input head just past the consumed cells.
\end{lemma}

\begin{lemma}[Completed parser run]
\lean{Turing.timedPrefix_complete}
\label{Turing.timedPrefix_complete}
\leanok
\uses{Turing.MultiTapeTM.initCfg, Turing.MultiTapeTM.runFrom,
  Turing.MultiTapeTM.runFrom_succ_eq_step', Turing.pairEncode,
  Turing.timedPrefixCfg, Turing.timedPrefixTM, Turing.timedPrefix_advance,
  Turing.timedPrefix_code, Turing.timed_code_get, Turing.timed_input_length,
  Turing.timed_pair_separator}
\difficulty{4}
Running the prefix parser $4|b| + 2|\alpha| + 6$ steps from its initial
configuration on the nested pair encoding
$\langle \langle b, \alpha\rangle, x\rangle$ halts it with the full clock $b$
stored on the work tape, output exactly $\alpha$, and the input head parked at
the start of the suffix region holding $x$; both delimiters are consumed, also
when $b$ or $\alpha$ is empty.
\end{lemma}

\begin{definition}[Value of a clock word]
\lean{Turing.timedValue}
\label{Turing.timedValue}
\leanok
The little-endian numeric value of a fixed-width Boolean word: the empty word
has value $0$, and the word with first bit $b$ and tail $w$ has value
$2 \cdot \mathrm{value}(w) + b$.
\end{definition}

\begin{lemma}[Canonical clock words]
\lean{Turing.timedValue_bits}
\label{Turing.timedValue_bits}
\leanok
\uses{Turing.timedValue}
\difficulty{3}
For every natural number $t$, the little-endian value of the binary
representation of $t$ is $t$.
\end{lemma}

\begin{definition}[Fixed-width borrow subtraction]
\lean{Turing.timedBorrow}
\label{Turing.timedBorrow}
\leanok
Fixed-width binary subtraction of an incoming borrow flag from a little-endian
word: each output bit is the exclusive or of the corresponding input bit with
the propagated borrow, the borrow propagates exactly through the zero bits,
and the result is the final (underflow) borrow flag together with a word of
the same width.
\end{definition}

\begin{lemma}[A cleared borrow is the identity]
\lean{Turing.timedBorrow_false}
\label{Turing.timedBorrow_false}
\leanok
\uses{Turing.timedBorrow}
\difficulty{3}
Subtracting a cleared borrow flag from a word leaves the word unchanged and
reports no underflow.
\end{lemma}

\begin{lemma}[Borrow subtraction preserves width]
\lean{Turing.timedBorrow_length}
\label{Turing.timedBorrow_length}
\leanok
\uses{Turing.timedBorrow}
\difficulty{3}
For every incoming borrow flag and word, the word returned by the fixed-width
borrow subtraction has the same length as the input word.
\end{lemma}

\begin{lemma}[Underflow detects a zero budget]
\lean{Turing.timedBorrow_underflow}
\label{Turing.timedBorrow_underflow}
\leanok
\uses{Turing.timedBorrow, Turing.timedBorrow_false, Turing.timedValue}
\difficulty{3}
Subtracting a set borrow flag from a word underflows if and only if the word's
little-endian value is $0$.
\end{lemma}

\begin{lemma}[A successful borrow decrements the value]
\lean{Turing.timedBorrow_value}
\label{Turing.timedBorrow_value}
\leanok
\uses{Turing.TimedControl, Turing.timedBorrow, Turing.timedBorrow_false,
  Turing.timedValue}
\difficulty{4}
If a word has positive little-endian value, then the value of the word
obtained by subtracting a set borrow flag is exactly one less than the value
of the original word.
\end{lemma}

\begin{definition}[Timed-interpreter control states]
\lean{Turing.TimedControl}
\label{Turing.TimedControl}
\leanok
\uses{Turing.UniversalControl}
The finite control of the timed interpreter: a working copy of the table
interpreter's control, clock phases (rewind and borrow) that retain a selected
action's eight decoded bits and halting flag while the countdown is serviced,
an execute phase, and three output phases (emission start, buffer rewind, and
flush).
\end{definition}

\begin{definition}[Six-lane assembly]
\lean{Turing.timedSix}
\label{Turing.timedSix}
\leanok
The assembly of a function on six tape indices from four core values, a clock
value, and a buffer value: indices $0$ through $3$ map to the core values,
index $4$ to the clock value, and index $5$ to the buffer value.
\end{definition}

\begin{definition}[Lifting an interpreter action]
\lean{Turing.timedAction}
\label{Turing.timedAction}
\leanok
\uses{Turing.Action, Turing.TimedControl, Turing.UniversalControl,
  Turing.timedSix}
The lift of a four-tape interpreter action to six tapes: the four core lanes
act as before, the clock lane is untouched, any emission is written to the
output-buffer lane (advancing its head) instead of being emitted, and a
halting successor is redirected to the emission-start phase so that the
machine stays live.
\end{definition}

\begin{definition}[Timed administrative action]
\lean{Turing.timedAdmin}
\label{Turing.timedAdmin}
\leanok
\uses{Turing.Action, Turing.TimedControl, Turing.timedSix}
A six-tape action that leaves the four core lanes and the input head
untouched, performs a designated optional write and move on the clock lane and
on the output-buffer lane, optionally emits one real output bit, and passes to
a designated control.
\end{definition}

\begin{definition}[The timed interpreter]
\lean{Turing.timedInterpreter}
\label{Turing.timedInterpreter}
\leanok
\uses{Turing.MultiTapeTM, Turing.TimedControl, Turing.timedAction,
  Turing.timedAdmin, Turing.universalInterpreter}
The six-tape interpreter with a countdown clock: it follows the table
interpreter, buffering all emissions, until a decoded action is pending; it
then rewinds the clock lane and performs one fixed-width borrow pass.  A
successful borrow executes the retained action; an underflow halts the machine
emitting only the timeout tag $0$.  When the simulated machine halts, the
machine emits the success tag $1$ and flushes the buffered output before
halting.
\end{definition}

\begin{definition}[Lifting a four-tape configuration]
\lean{Turing.timedLift}
\label{Turing.timedLift}
\leanok
\uses{Turing.Cfg, Turing.FinTM.bufferTape, Turing.TimedControl,
  Turing.UniversalControl, Turing.timedSix}
The six-tape configuration representing a four-tape interpreter configuration
together with a clock word: the four core lanes are copied, the clock lane
holds the clock word as a buffer with head at its right blank, the buffer lane
holds the four-tape configuration's accumulated output, and the real output is
empty; a halted four-tape control maps to the emission-start phase.
\end{definition}

\begin{lemma}[Non-record states are unaffected by stopping]
\lean{Turing.timedCut_regular}
\label{Turing.timedCut_regular}
\leanok
\uses{Turing.UniversalControl, Turing.timedAction_apply,
  Turing.timedCutInterpreter, Turing.universalInterpreter}
\difficulty{2}
For every control state that is not a pending-action state, the transition
function of the stopped table interpreter agrees with that of the original
table interpreter on all reads.
\end{lemma}

\begin{lemma}[Liveness propagates backwards]
\lean{Turing.timedCut_live_before}
\label{Turing.timedCut_live_before}
\leanok
\uses{Turing.Cfg, Turing.MultiTapeTM.runFrom, Turing.MultiTapeTM.runFrom_add,
  Turing.MultiTapeTM.runFrom_of_halt, Turing.UniversalControl,
  Turing.timedAction_apply, Turing.timedCutInterpreter}
\difficulty{3}
If the stopped table interpreter has not halted after $t$ steps from a
configuration, then it has not halted after any $s \le t$ steps from that
configuration (halting is absorbing).
\end{lemma}

\begin{lemma}[No record before a live endpoint]
\lean{Turing.timedCut_no_record}
\label{Turing.timedCut_no_record}
\leanok
\uses{Turing.Action.apply, Turing.Cfg, Turing.MultiTapeTM.runFrom,
  Turing.MultiTapeTM.runFrom_succ_eq_step', Turing.MultiTapeTM.step,
  Turing.UniversalControl, Turing.timedCutInterpreter,
  Turing.timedCut_live_before}
\difficulty{3}
If the stopped table interpreter has not halted after $t$ steps from a
configuration, then at every time $s < t$ its control is not a pending-action
state, because the stopped interpreter halts one step after reaching such a
state.
\end{lemma}

\begin{lemma}[Core reads of the six-lane assembly]
\lean{Turing.timedSix_core}
\label{Turing.timedSix_core}
\leanok
\uses{Turing.timedSix}
\difficulty{2}
Restricting the six-lane assembly of four core values, a clock value, and a
buffer value to the first four indices returns the four core values.
\end{lemma}

\begin{lemma}[Applying a lifted action]
\lean{Turing.timedAction_apply}
\label{Turing.timedAction_apply}
\leanok
\uses{Turing.Action, Turing.Action.apply, Turing.Cfg,
  Turing.FinTM.bufferTape_append, Turing.UniversalControl,
  Turing.timedAction, Turing.timedLift, Turing.timedSix}
\difficulty{3}
Applying the six-tape lift of a four-tape interpreter action to the six-tape
lift of a four-tape configuration (with any clock word) equals the lift of the
configuration obtained by applying the action: the clock word is unchanged and
any emission is captured on the buffer lane, including an emission on the
action's halting transition.
\end{lemma}

\begin{lemma}[One regular step replays]
\lean{Turing.timed_regular_step}
\label{Turing.timed_regular_step}
\leanok
\uses{Turing.Cfg, Turing.Cfg.inputSymbol, Turing.Cfg.workTapeSymbols,
  Turing.MultiTapeTM.step, Turing.UniversalControl, Turing.timedAction,
  Turing.timedAction_apply, Turing.timedCutInterpreter,
  Turing.timedCut_regular, Turing.timedInterpreter, Turing.timedLift,
  Turing.universalInterpreter}
\difficulty{4}
For every live four-tape configuration whose control is not a pending-action
state, one step of the timed interpreter from its six-tape lift (with any
clock word) equals the lift of one step of the stopped table interpreter, with
the same clock word.
\end{lemma}

\begin{lemma}[Replaying a stopped run]
\lean{Turing.timed_replay}
\label{Turing.timed_replay}
\leanok
\uses{Turing.Cfg, Turing.MultiTapeTM.runFrom,
  Turing.MultiTapeTM.runFrom_succ_eq_step', Turing.UniversalControl,
  Turing.timedCutInterpreter, Turing.timedCut_live_before,
  Turing.timedCut_no_record, Turing.timedInterpreter, Turing.timedLift,
  Turing.timed_regular_step}
\difficulty{4}
If the stopped table interpreter has not halted after $t$ steps from a
four-tape configuration, then $t$ steps of the timed interpreter from the
six-tape lift of that configuration equal the lift of the $t$-step stopped
run, with an unchanged clock word; no countdown or output phase is visited in
the interior of the run.
\end{lemma}

\begin{definition}[Single active lane configuration]
\lean{Turing.timed_laneCfg}
\label{Turing.timed_laneCfg}
\leanok
\uses{Turing.Cfg, Turing.FinTM.sweepTape}
A configuration in which one designated work tape holds a finite zipper --- a
frontier position, a stack of processed cells, and a list of unprocessed cells
--- while every other component (the remaining tapes, all other heads, the
input position, and the output) is taken from a base configuration.
\end{definition}

\begin{definition}[Single-lane action]
\lean{Turing.timed_laneAction}
\label{Turing.timed_laneAction}
\leanok
\uses{Turing.Action}
An action that writes an optional symbol and moves on exactly one designated
work tape, leaves the input head, the output, and all other work tapes
stationary, and passes to a given control state.
\end{definition}

\begin{lemma}[Reading the active lane]
\lean{Turing.timed_laneCfg_read}
\label{Turing.timed_laneCfg_read}
\leanok
\uses{Turing.Cfg, Turing.Cfg.workTapeSymbols, Turing.FinTM.sweepTape_read,
  Turing.timed_laneCfg}
\difficulty{1}
In a single-active-lane configuration, the active lane reads the first entry
of the unprocessed zipper list, or a blank when that list is empty.
\end{lemma}

\begin{lemma}[Right-moving zipper step on one lane]
\lean{Turing.timed_laneCfg_right}
\label{Turing.timed_laneCfg_right}
\leanok
\uses{Turing.Action.apply, Turing.Cfg, Turing.FinTM.sweepTape_right,
  Turing.moveInputPos_zero, Turing.timed_laneAction, Turing.timed_laneCfg}
\difficulty{3}
Applying a right-moving single-lane action that writes a symbol $b$ to a
single-active-lane configuration whose unprocessed list begins with a cell $a$
consumes $a$, pushes $b$ onto the processed stack, advances the frontier by
one, and passes to the designated control; all inactive data is unchanged.
\end{lemma}

\begin{lemma}[Generic one-lane transduction]
\lean{Turing.timed_lane_run}
\label{Turing.timed_lane_run}
\leanok
\uses{Turing.Action.apply, Turing.Cfg, Turing.Cfg.workTapeSymbols,
  Turing.FinTM.sweepFold, Turing.MultiTapeTM, Turing.MultiTapeTM.runFrom,
  Turing.MultiTapeTM.runFrom_succ_eq_step, Turing.MultiTapeTM.runFrom_zero,
  Turing.timed_laneAction, Turing.timed_laneCfg, Turing.timed_laneCfg_read,
  Turing.timed_laneCfg_right}
\difficulty{5}
Suppose a machine's transitions, on a family of control states indexed by a
register type, always perform a right-moving single-lane write determined by a
local visit function of the register and the symbol read on the designated
lane.  Then running the machine for exactly the length of a word laid out on
that lane's zipper computes the left fold of the visit function over the word:
the final control carries the folded register, the frontier advances by the
word's length, the processed stack holds the rewritten word reversed, and all
inactive tapes are unchanged.
\end{lemma}

\begin{lemma}[Clock writes are single-lane actions]
\lean{Turing.timedAdmin_clock}
\label{Turing.timedAdmin_clock}
\leanok
\uses{Turing.TimedControl, Turing.timedAdmin, Turing.timed_laneAction}
\difficulty{2}
A timed administrative action that writes and moves only on the clock lane
coincides with the generic single-lane action on lane $4$ of the six-tape
machine.
\end{lemma}

\begin{lemma}[Borrow as a local fold]
\lean{Turing.timedBorrow_fold}
\label{Turing.timedBorrow_fold}
\leanok
\uses{Turing.FinTM.sweepFold, Turing.timedBorrow}
\difficulty{3}
The fixed-width borrow subtraction equals the left fold of the local rule
sending a carry flag $c$ and a bit $b$ to the new carry $c \wedge \lnot b$ and
the output bit $b \oplus c$.
\end{lemma}

\begin{lemma}[Exact borrow transduction]
\lean{Turing.timed_borrow_run}
\label{Turing.timed_borrow_run}
\leanok
\uses{Turing.Cfg, Turing.MultiTapeTM.runFrom, Turing.TimedControl,
  Turing.timedAdmin_clock, Turing.timedBorrow, Turing.timedBorrow_fold,
  Turing.timedInterpreter, Turing.timed_laneCfg, Turing.timed_lane_run}
\difficulty{3}
From the borrow phase with a given carry flag and the clock word laid out on
the clock lane's zipper, the timed interpreter processes the word in exactly
its length: the final control carries the resulting underflow flag, the
processed stack holds the borrowed word, and neither the four simulation lanes
nor the buffered output is touched.
\end{lemma}

\begin{lemma}[Shifting a zipper frontier]
\lean{Turing.timed_sweep_shift}
\label{Turing.timed_sweep_shift}
\leanok
\uses{Turing.FinTM.sweepTape, Turing.FinTM.sweepTape_read,
  Turing.FinTM.sweepTape_right}
\difficulty{4}
Moving the frontier of a finite zipper across a segment $w$ --- pushing the
reversal of $w$ onto the processed stack and advancing the frontier by $|w|$
--- leaves the represented tape unchanged.
\end{lemma}

\begin{lemma}[A buffer is a zipper]
\lean{Turing.timed_buffer_zipper}
\label{Turing.timed_buffer_zipper}
\leanok
\uses{Turing.FinTM.bufferTape, Turing.FinTM.sweepTape}
\difficulty{3}
The buffer tape of a word equals the zipper with frontier $0$, empty processed
stack, and the word as the unprocessed list.
\end{lemma}

\begin{lemma}[A buffer as a fully processed zipper]
\lean{Turing.timed_buffer_zipper_end}
\label{Turing.timed_buffer_zipper_end}
\leanok
\uses{Turing.FinTM.bufferTape, Turing.FinTM.sweepTape,
  Turing.timed_buffer_zipper, Turing.timed_sweep_shift}
\difficulty{2}
The buffer tape of a word also equals the zipper with frontier at the word's
length, the reversed word as the processed stack, and an empty unprocessed
list.
\end{lemma}

\begin{definition}[Clock-phase configuration]
\lean{Turing.timedClockCfg}
\label{Turing.timedClockCfg}
\leanok
\uses{Turing.Cfg, Turing.FinTM.bufferTape, Turing.TimedControl}
The configuration obtained from a base six-tape configuration by overriding
only the finite control and the clock lane, whose contents become the buffer
of a given word with the head at a given position.
\end{definition}

\begin{lemma}[One clock-phase step]
\lean{Turing.timedClock_step}
\label{Turing.timedClock_step}
\leanok
\uses{Turing.Action.apply, Turing.Cfg, Turing.Cfg.workTapeSymbols,
  Turing.FinTM.bufferTape, Turing.MultiTapeTM.step, Turing.TimedControl,
  Turing.moveInputPos_zero, Turing.timedAdmin, Turing.timedClockCfg,
  Turing.timedInterpreter, Turing.timedSix}
\difficulty{3}
Suppose that, whenever the clock lane reads its designated symbol, the timed
interpreter's transition in control $q$ is a pure clock move by
$d \in \{-1,0,1\}$ (no writes, no emission) into a control $q'$.  Then one
step from the clock-phase configuration moves only the clock head by $d$ and
enters $q'$.
\end{lemma}

\begin{lemma}[Rewinding the clock]
\lean{Turing.timed_clock_back}
\label{Turing.timed_clock_back}
\leanok
\uses{Turing.Cfg, Turing.FinTM.bufferTape, Turing.FinTM.bufferTape_nat,
  Turing.MultiTapeTM.runFrom, Turing.MultiTapeTM.runFrom_succ_eq_step,
  Turing.MultiTapeTM.runFrom_zero, Turing.TimedControl,
  Turing.timedClockCfg, Turing.timedClock_step, Turing.timedInterpreter}
\difficulty{4}
For every $j$ at most the clock word's length, from the clock-rewind phase
with the clock head at $j-1$, exactly $j+1$ steps of the timed interpreter
reach the borrow phase with the borrow flag set and the clock head at $0$; the
retained action bits and halting flag are unchanged.
\end{lemma}

\begin{lemma}[One borrow pass over the clock]
\lean{Turing.timed_clock_borrow}
\label{Turing.timed_clock_borrow}
\leanok
\uses{Turing.Cfg, Turing.MultiTapeTM.runFrom, Turing.TimedControl,
  Turing.timedBorrow, Turing.timedBorrow_length, Turing.timedClockCfg,
  Turing.timedInterpreter, Turing.timed_borrow_run,
  Turing.timed_buffer_zipper, Turing.timed_buffer_zipper_end,
  Turing.timed_laneCfg}
\difficulty{4}
From the borrow phase with a given carry flag and the clock head at $0$,
exactly as many steps of the timed interpreter as the clock word's length
replace the word by its borrow subtraction, ending at the right blank with the
resulting underflow flag in the control; the word's width is preserved.
\end{lemma}

\begin{lemma}[Entering the clock rewind]
\lean{Turing.timed_clock_start}
\label{Turing.timed_clock_start}
\leanok
\uses{Turing.Action.apply, Turing.Cfg, Turing.MultiTapeTM.step,
  Turing.UniversalControl, Turing.moveInputPos_zero, Turing.timedAdmin,
  Turing.timedClockCfg, Turing.timedInterpreter, Turing.timedLift,
  Turing.timedSix}
\difficulty{3}
From the six-tape lift of a four-tape configuration whose control is a
pending-action state, one step of the timed interpreter enters the
clock-rewind phase --- retaining the action's decoded bits and halting flag
--- with the clock head moved to the last clock bit, and without applying any
source transition.
\end{lemma}

\begin{lemma}[The complete countdown pass]
\lean{Turing.timed_clock_pass}
\label{Turing.timed_clock_pass}
\leanok
\uses{Turing.Cfg, Turing.MultiTapeTM.runFrom, Turing.UniversalControl,
  Turing.timedBorrow, Turing.timedClockCfg, Turing.timedCut_run_join,
  Turing.timedInterpreter, Turing.timedLift, Turing.timed_clock_back,
  Turing.timed_clock_borrow, Turing.timed_clock_start}
\difficulty{4}
From the six-tape lift of a pending-action configuration with clock word $w$,
exactly $2|w| + 2$ steps of the timed interpreter complete the clock rewind
and the borrow pass, reaching the borrow phase at the right blank with the
subtracted word and its underflow flag in the control.
\end{lemma}

\begin{lemma}[Executing the retained action]
\lean{Turing.timed_execute}
\label{Turing.timed_execute}
\leanok
\uses{Turing.Action.apply, Turing.Cfg, Turing.Cfg.inputSymbol,
  Turing.Cfg.workTapeSymbols, Turing.MultiTapeTM.step,
  Turing.UniversalControl, Turing.timedAction, Turing.timedAction_apply,
  Turing.timedInterpreter, Turing.timedLift, Turing.universalInterpreter}
\difficulty{3}
From the six-tape lift of a pending-action configuration, with the control
replaced by the execute phase, one step of the timed interpreter applies the
retained source action: the result equals the lift of one step of the
unstopped table interpreter from the four-tape configuration, with the clock
word unchanged.
\end{lemma}

\begin{lemma}[A positive budget executes one transition]
\lean{Turing.timed_clock_success}
\label{Turing.timed_clock_success}
\leanok
\uses{Turing.Action.apply, Turing.Cfg, Turing.Cfg.inputSymbol,
  Turing.Cfg.workTapeSymbols, Turing.FinTM.bufferTape_nat,
  Turing.MultiTapeTM.runFrom, Turing.MultiTapeTM.runFrom_succ_eq_step',
  Turing.MultiTapeTM.step, Turing.UniversalControl,
  Turing.moveInputPos_zero, Turing.timedAdmin, Turing.timedBorrow,
  Turing.timedBorrow_length, Turing.timedBorrow_underflow,
  Turing.timedClockCfg, Turing.timedInterpreter, Turing.timedLift,
  Turing.timedSix, Turing.timedValue, Turing.timed_clock_pass,
  Turing.timed_execute, Turing.universalInterpreter}
\difficulty{5}
If the clock word $w$ has positive little-endian value, then from the six-tape
lift of a pending-action configuration, exactly $2|w| + 4$ steps of the timed
interpreter decrement the budget once and apply the retained action: the
result is the lift of one step of the unstopped table interpreter, with clock
word the borrow subtraction of $w$.
\end{lemma}

\begin{lemma}[A zero budget times out]
\lean{Turing.timed_clock_timeout}
\label{Turing.timed_clock_timeout}
\leanok
\uses{Turing.Action.apply, Turing.Cfg, Turing.Cfg.inputSymbol,
  Turing.Cfg.workTapeSymbols, Turing.FinTM.bufferTape_nat,
  Turing.MultiTapeTM.runFrom, Turing.MultiTapeTM.runFrom_succ_eq_step',
  Turing.MultiTapeTM.step, Turing.UniversalControl, Turing.timedAdmin,
  Turing.timedBorrow, Turing.timedBorrow_length,
  Turing.timedBorrow_underflow, Turing.timedClockCfg,
  Turing.timedInterpreter, Turing.timedLift, Turing.timedValue,
  Turing.timed_clock_pass}
\difficulty{4}
If the clock word $w$ has little-endian value $0$, then after $2|w| + 3$ steps
of the timed interpreter from the six-tape lift of a pending-action
configuration, the machine has halted with real output exactly the single
timeout tag $0$, regardless of any emissions buffered so far.
\end{lemma}

\begin{definition}[Output-phase configuration]
\lean{Turing.timedOutputCfg}
\label{Turing.timedOutputCfg}
\leanok
\uses{Turing.Cfg, Turing.TimedControl, Turing.UniversalControl,
  Turing.timedLift}
The configuration obtained from the six-tape lift of a four-tape configuration
and a clock word by overriding the finite control, the output-buffer head
position, and the real output; all simulation and clock tapes are retained.
\end{definition}

\begin{lemma}[One output-phase step]
\lean{Turing.timed_output_step}
\label{Turing.timed_output_step}
\leanok
\uses{Turing.Action.apply, Turing.Cfg, Turing.Cfg.workTapeSymbols,
  Turing.FinTM.bufferTape, Turing.MultiTapeTM.step, Turing.TimedControl,
  Turing.UniversalControl, Turing.moveInputPos_zero, Turing.timedAdmin,
  Turing.timedInterpreter, Turing.timedLift, Turing.timedOutputCfg,
  Turing.timedSix}
\difficulty{3}
Suppose that, whenever the output-buffer lane reads its designated symbol, the
timed interpreter's transition in control $q$ is a pure buffer move by
$d \in \{-1,0,1\}$ with an optional real emission.  Then one step from the
output-phase configuration moves only the buffer head by $d$, appends the
optional emission to the real output, and passes to the designated control.
\end{lemma}

\begin{lemma}[Rewinding the output buffer]
\lean{Turing.timed_output_back}
\label{Turing.timed_output_back}
\leanok
\uses{Turing.Cfg, Turing.FinTM.bufferTape, Turing.FinTM.bufferTape_nat,
  Turing.MultiTapeTM.runFrom, Turing.MultiTapeTM.runFrom_succ_eq_step,
  Turing.MultiTapeTM.runFrom_zero, Turing.UniversalControl,
  Turing.timedInterpreter, Turing.timedOutputCfg, Turing.timed_output_step}
\difficulty{4}
For every $j$ at most the buffered output's length, from the buffer-rewind
phase with the buffer head at $j-1$, exactly $j+1$ steps of the timed
interpreter reach the flush phase with the buffer head at $0$, emitting the
success tag $1$ at the left blank before any data bit.
\end{lemma}

\begin{lemma}[Flushing the buffer]
\lean{Turing.timed_output_forward}
\label{Turing.timed_output_forward}
\leanok
\uses{Turing.Cfg, Turing.FinTM.bufferTape, Turing.FinTM.bufferTape_nat,
  Turing.MultiTapeTM.runFrom, Turing.MultiTapeTM.runFrom_succ_eq_step,
  Turing.UniversalControl, Turing.timedInterpreter, Turing.timedOutputCfg,
  Turing.timed_output_step, Turing.universal_table_read}
\difficulty{4}
If the buffered output splits as a processed prefix $l$ followed by a
remainder $r$, with the buffer head at $|l|$, then the flush phase of the
timed interpreter emits exactly the bits of $r$, in order, in $|r| + 1$ steps
and halts at the right blank without a further emission.
\end{lemma}

\begin{lemma}[Halting flushes tag and buffer]
\lean{Turing.timed_flush}
\label{Turing.timed_flush}
\leanok
\uses{Turing.Cfg, Turing.MultiTapeTM.runFrom, Turing.UniversalControl,
  Turing.timedCut_run_join, Turing.timedInterpreter, Turing.timedLift,
  Turing.timedOutputCfg, Turing.timed_output_back,
  Turing.timed_output_forward, Turing.timed_output_step}
\difficulty{4}
From the six-tape lift of a halted four-tape configuration whose buffered
output is a word $w$, exactly $2|w| + 3$ steps of the timed interpreter halt
the machine with real output the success tag $1$ followed by exactly $w$.
\end{lemma}

\begin{lemma}[The parser is live at the penultimate step]
\lean{Turing.timedPrefix_penultimate}
\label{Turing.timedPrefix_penultimate}
\leanok
\uses{Turing.MultiTapeTM.initCfg, Turing.MultiTapeTM.runFrom,
  Turing.MultiTapeTM.runFrom_succ_eq_step', Turing.pairEncode,
  Turing.timedPrefixTM, Turing.timedPrefix_advance, Turing.timedPrefix_code,
  Turing.timed_code_get, Turing.timed_input_length,
  Turing.timed_pair_separator}
\difficulty{3}
After $4|b| + 2|\alpha| + 5$ steps on the nested pair encoding
$\langle \langle b, \alpha\rangle, x\rangle$ --- one step before consuming
the final separator cell --- the prefix parser has not halted.
\end{lemma}

\begin{lemma}[The parser is live before completion]
\lean{Turing.timedPrefix_live}
\label{Turing.timedPrefix_live}
\leanok
\uses{Turing.MultiTapeTM.initCfg, Turing.MultiTapeTM.runFrom,
  Turing.MultiTapeTM.runFrom_add, Turing.MultiTapeTM.runFrom_of_halt,
  Turing.pairEncode, Turing.timedPrefixTM, Turing.timedPrefix_penultimate}
\difficulty{3}
At every time strictly before $4|b| + 2|\alpha| + 6$ steps on the nested pair
encoding $\langle \langle b, \alpha\rangle, x\rangle$, the prefix parser has
not halted.
\end{lemma}

\begin{definition}[The composed clock-preserving canonizer]
\lean{Turing.timedCanonTM}
\label{Turing.timedCanonTM}
\leanok
\uses{Turing.EffectiveMachineCode, Turing.FinTM, Turing.FinTM.bufferedCompTM,
  Turing.timedPrefixTM}
The buffered composition of the prefix parser with the coding scheme's
canonizer: the parser extracts the code string from the nested pair encoding,
and only that extracted code is supplied as virtual input to the canonizer.
\end{definition}

\begin{definition}[The retained clock lane]
\lean{Turing.timedCanonClock}
\label{Turing.timedCanonClock}
\leanok
\uses{Turing.EffectiveMachineCode, Turing.timedCanonTM}
The tape index of the composed clock-preserving canonizer corresponding to the
prefix parser's unique work tape, on which the extracted clock word is
retained.
\end{definition}

\begin{lemma}[Entering the canonizer]
\lean{Turing.timedCanon_start}
\label{Turing.timedCanon_start}
\leanok
\uses{Turing.EffectiveMachineCode, Turing.FinTM.bufferTape,
  Turing.FinTM.bufferedCompTM, Turing.FinTM.bufferedFirstCfg_init,
  Turing.FinTM.bufferedFirstCfg_rewind, Turing.FinTM.bufferedFirstCfg_run,
  Turing.FinTM.bufferedSecondCfg, Turing.MultiTapeTM.initCfg,
  Turing.MultiTapeTM.runFrom, Turing.MultiTapeTM.runFrom_add,
  Turing.pairEncode, Turing.timedCanonTM, Turing.timedPrefixTM,
  Turing.timedPrefix_complete, Turing.timedPrefix_live,
  Turing.timed_input_length}
\difficulty{3}
On the nested pair encoding $\langle \langle b, \alpha\rangle, x\rangle$,
exactly $4|b| + 3|\alpha| + 8$ steps of the composed clock-preserving
canonizer complete the parse and hand the extracted code $\alpha$ to the
canonizer in its initial configuration, with the clock word $b$ preserved on
its lane (head at $|b|$) and the physical input head parked at the start of
the suffix region.
\end{lemma}

\begin{lemma}[Running the canonizer in place]
\lean{Turing.timedCanon_run}
\label{Turing.timedCanon_run}
\leanok
\uses{Turing.EffectiveMachineCode, Turing.FinTM.VirtualTag,
  Turing.FinTM.bufferTape, Turing.FinTM.bufferedSecondCfg,
  Turing.FinTM.bufferedSecondCfg_run, Turing.MultiTapeTM.initCfg,
  Turing.MultiTapeTM.runFrom, Turing.MultiTapeTM.runFrom_add,
  Turing.pairEncode, Turing.timedCanonTM, Turing.timedCanon_start,
  Turing.timedPrefixTM, Turing.timed_input_length}
\difficulty{3}
For every $t$, running the composed clock-preserving canonizer
$4|b| + 3|\alpha| + 8 + t$ steps on the nested pair encoding equals, for some
internal rewind flag, the buffered second-phase configuration of the canonizer
run for $t$ steps on virtual input $\alpha$; the physical input head and the
retained clock word $b$ remain stationary.
\end{lemma}

\begin{lemma}[Canonizer completion]
\lean{Turing.timedCanon_complete}
\label{Turing.timedCanon_complete}
\leanok
\uses{Turing.CodeTM.serialize, Turing.EffectiveMachineCode,
  Turing.FinTM.bufferTape, Turing.FinTM.bufferedSecondCfg,
  Turing.FinTM.computesInTime_iff, Turing.MultiTapeTM.initCfg,
  Turing.MultiTapeTM.runFrom, Turing.pairEncode, Turing.timedCanonClock,
  Turing.timedCanonTM, Turing.timedCanon_run}
\difficulty{3}
After $4|b| + 3|\alpha| + 8$ plus the canonizer's time bound at $|\alpha|$
steps on the nested pair encoding, the composed clock-preserving canonizer has
halted; its output is the serialization of the machine decoded from $\alpha$,
its input head is parked at position $4|b| + 2|\alpha| + 7$, and its retained
clock lane still holds the buffer of $b$ with head at $|b|$.
\end{lemma}

\begin{definition}[Deadline-inclusive answer]
\lean{Turing.timedAnswer}
\label{Turing.timedAnswer}
\leanok
\uses{Turing.Cfg, Turing.CodeTM, Turing.MultiTapeTM.runFrom}
The answer of a source configuration of a coded machine after a deadline of
$t$ steps: the success tag $1$ followed by the accumulated output when the
machine has halted within $t$ steps, and the single timeout tag $0$ otherwise.
\end{definition}

\begin{lemma}[The timed interpreter finishes uniformly]
\lean{Turing.timed_interpret_finishes}
\label{Turing.timed_interpret_finishes}
\leanok
\uses{Turing.Cfg, Turing.CodeTM, Turing.CodeTM.serialize,
  Turing.MultiTapeTM.outputSymbol, Turing.MultiTapeTM.runFrom,
  Turing.MultiTapeTM.runFrom_add, Turing.MultiTapeTM.runFrom_of_halt,
  Turing.MultiTapeTM.runFrom_succ_eq_step, Turing.MultiTapeTM.runFrom_zero,
  Turing.MultiTapeTM.step, Turing.MultiTapeTM.step_output,
  Turing.timedAnswer, Turing.timedBorrow, Turing.timedBorrow_length,
  Turing.timedBorrow_value, Turing.timedCut_live_block,
  Turing.timedInterpreter, Turing.timedLift, Turing.timedValue,
  Turing.timed_clock_success, Turing.timed_clock_timeout, Turing.timed_flush,
  Turing.timed_replay, Turing.universalSimulationCfg}
\difficulty{6}
Let $M$ be a coded machine with serialization of length $L$ and state count
$N+1$, let $\mathrm{src}$ be any configuration of $M$, let the table cursor be
at most $L$, and let the clock word have width $w$ and little-endian value
$r$.  Then the timed interpreter, started from the six-tape lift of the
interpreter checkpoint of $\mathrm{src}$ with that clock word, halts within
\[
  \bigl(3L + 5(N+1) + 20 + 2w + 8\bigr)\,(r+1)
    + 2\,\lvert \text{output of } \mathrm{src}\rvert
\]
steps, and its real output is exactly the deadline-inclusive answer of
$\mathrm{src}$ at budget $r$: halting on the final allowed transition counts
as success.
\end{lemma}

\begin{lemma}[Width of the binary clock]
\lean{Turing.timed_bits_length}
\label{Turing.timed_bits_length}
\leanok
\difficulty{3}
For every natural number $t$, the binary representation of $t$ has at most $t$
bits; in particular the bound also holds at $t = 0$.
\end{lemma}

\begin{definition}[Five-lane assembly]
\lean{Turing.timedFive}
\label{Turing.timedFive}
\leanok
The assembly of a function on five tape indices from four core values and a
buffer value: indices $0$ through $3$ map to the core values (table, state,
simulated work, and input marker) and index $4$ to the buffer value.
\end{definition}

\begin{definition}[Framed interpreter action]
\lean{Turing.timedFrameAction}
\label{Turing.timedFrameAction}
\leanok
\uses{Turing.Action, Turing.FinTM, Turing.TimedControl, Turing.timedFive}
The lift of a six-tape timed-interpreter action to the assembled machine with
$k + 5$ tapes, for a base machine with $k$ tapes and a designated clock tape:
the action's clock component acts on the retained clock lane of the old block,
its remaining five components act on the five fresh lanes, and all other old
lanes are left untouched.
\end{definition}

\begin{definition}[The capture wrapper]
\lean{Turing.timedCaptureTM}
\label{Turing.timedCaptureTM}
\leanok
\uses{Turing.FinTM, Turing.FinTM.controlAction, Turing.FinTM.tapeBlocks,
  Turing.TimedControl, Turing.timedFrameAction, Turing.timedInterpreter,
  Turing.timedSix}
The capture wrapper of a machine $M$ with a designated clock tape: it runs $M$
in lockstep while writing $M$'s emissions to a fresh table lane instead of the
output; when $M$ halts, control transfers to the timed interpreter, which runs
on the captured table, the retained clock lane, and four further fresh lanes.
\end{definition}

\begin{definition}[First-phase capture configuration]
\lean{Turing.timedCaptureCfg}
\label{Turing.timedCaptureCfg}
\leanok
\uses{Turing.Cfg, Turing.FinTM, Turing.FinTM.bufferTape,
  Turing.FinTM.tapeBlocks, Turing.timedCaptureTM}
The first-phase configuration of the capture wrapper corresponding to a
configuration of the base machine $M$: $M$'s tapes and heads are embedded, the
captured-table lane holds $M$'s accumulated output as a buffer with head at
its right blank, the remaining fresh lanes are blank, and the real output is
empty.
\end{definition}

\begin{lemma}[Capture starts blank]
\lean{Turing.timedCapture_init}
\label{Turing.timedCapture_init}
\leanok
\uses{Turing.FinTM, Turing.FinTM.tapeBlocks, Turing.MultiTapeTM.initCfg,
  Turing.timedCaptureCfg, Turing.timedCaptureTM}
\difficulty{3}
On every input, the initial configuration of the capture wrapper equals the
first-phase embedding of the base machine's initial configuration, with a
blank captured table and blank interpreter lanes.
\end{lemma}

\begin{lemma}[One captured step]
\lean{Turing.timedCapture_step}
\label{Turing.timedCapture_step}
\leanok
\uses{Turing.Action, Turing.Action.apply, Turing.Cfg, Turing.Cfg.inputSymbol,
  Turing.Cfg.workTapeSymbols, Turing.FinTM, Turing.FinTM.bufferTape_append,
  Turing.FinTM.tapeBlocks, Turing.MultiTapeTM.step, Turing.timedCaptureCfg,
  Turing.timedCaptureTM, Turing.timedFrame}
\difficulty{5}
For every live configuration of the base machine, one step of the capture
wrapper from its first-phase embedding equals the embedding of one step of the
base machine: the input head and the original work tapes move in lockstep, any
emission --- including one on the base machine's halting transition --- is
appended to the captured table rather than the real output, and the later
fresh lanes stay untouched; administrative states remain live.
\end{lemma}

\begin{lemma}[Lockstep capture of a run]
\lean{Turing.timedCapture_run}
\label{Turing.timedCapture_run}
\leanok
\uses{Turing.Cfg, Turing.FinTM, Turing.MultiTapeTM.runFrom,
  Turing.MultiTapeTM.runFrom_succ_eq_step', Turing.timedCaptureCfg,
  Turing.timedCaptureTM, Turing.timedCapture_step}
\difficulty{3}
If the base machine remains live at every time strictly below $t$, then $t$
steps of the capture wrapper from an embedded configuration equal the
embedding of $t$ steps of the base machine.
\end{lemma}

\begin{definition}[Interpreter-entry capture configuration]
\lean{Turing.timedCapturedCfg}
\label{Turing.timedCapturedCfg}
\leanok
\uses{Turing.Cfg, Turing.FinTM, Turing.timedCaptureCfg,
  Turing.timedCaptureTM, Turing.timedInterpreter}
The interpreter-entry configuration of the capture wrapper for a halted
configuration of the base machine: identical to its first-phase embedding
except that the finite control is the timed interpreter's initial state; the
retained clock lane becomes active again during interpretation.
\end{definition}

\begin{lemma}[Transferring to the interpreter]
\lean{Turing.timedCapture_transfer}
\label{Turing.timedCapture_transfer}
\leanok
\uses{Turing.Cfg, Turing.FinTM, Turing.MultiTapeTM.step,
  Turing.moveInputPos_zero, Turing.timedCaptureCfg, Turing.timedCaptureTM,
  Turing.timedCapturedCfg}
\difficulty{2}
From the first-phase embedding of a halted configuration of the base machine,
one step of the capture wrapper transfers control to the interpreter entry,
changing no tape contents, head positions, or output.
\end{lemma}

\begin{lemma}[Reaching the interpreter after capture]
\lean{Turing.timedCapture_start}
\label{Turing.timedCapture_start}
\leanok
\uses{Turing.FinTM, Turing.MultiTapeTM.initCfg, Turing.MultiTapeTM.runFrom,
  Turing.MultiTapeTM.runFrom_add, Turing.MultiTapeTM.runFrom_of_halt,
  Turing.MultiTapeTM.runFrom_succ_eq_step', Turing.timedCaptureTM,
  Turing.timedCapture_init, Turing.timedCapture_run,
  Turing.timedCapture_transfer, Turing.timedCapturedCfg}
\difficulty{4}
If the base machine halts on input $x$ within $T$ steps, then for some
$t \le T + 1$, running the capture wrapper $t$ steps from its initial
configuration on $x$ reaches the interpreter-entry configuration for the base
machine's configuration at time $T$; since halting is absorbing, all fields of
that configuration --- including the parked input head --- are retained, not
merely the completed output.
\end{lemma}

\begin{definition}[Framed interpreter configuration]
\lean{Turing.timedFrame}
\label{Turing.timedFrame}
\leanok
\uses{Turing.Cfg, Turing.FinTM, Turing.TimedControl, Turing.timedCaptureTM,
  Turing.timedFive}
The configuration of the assembled capture machine determined by a six-tape
timed-interpreter configuration together with inactive tapes and heads for the
old block: the interpreter's clock lane is placed at the retained clock index
of the old block, its remaining five lanes occupy the fresh indices, and every
other old lane carries the given inactive data.
\end{definition}

\begin{lemma}[Reads of a framed configuration]
\lean{Turing.timedFrame_reads}
\label{Turing.timedFrame_reads}
\leanok
\uses{Turing.Cfg, Turing.Cfg.workTapeSymbols, Turing.FinTM,
  Turing.TimedControl, Turing.timedFive, Turing.timedFrame, Turing.timedSix}
\difficulty{2}
The six symbols read by a framed configuration on the interpreter's active
lanes (the four core lanes, the retained clock lane, and the buffer lane)
coincide with the reads of the underlying six-tape interpreter configuration.
\end{lemma}

\begin{lemma}[Applying a framed action]
\lean{Turing.timedFrame_apply}
\label{Turing.timedFrame_apply}
\leanok
\uses{Turing.Action, Turing.Action.apply, Turing.Cfg, Turing.FinTM,
  Turing.TimedControl, Turing.timedFive, Turing.timedFrame,
  Turing.timedFrameAction}
\difficulty{4}
Applying the framed lift of a six-tape action to a framed configuration equals
the frame of the six-tape configuration obtained by applying the action; the
inactive lanes of the old block are unchanged.
\end{lemma}

\begin{lemma}[One framed step]
\lean{Turing.timedFrame_step}
\label{Turing.timedFrame_step}
\leanok
\uses{Turing.Cfg, Turing.Cfg.inputSymbol, Turing.FinTM,
  Turing.MultiTapeTM.step, Turing.MultiTapeTM.step_of_halt,
  Turing.TimedControl, Turing.timedCaptureTM, Turing.timedFrame,
  Turing.timedFrame_apply, Turing.timedFrame_reads, Turing.timedInterpreter}
\difficulty{3}
One step of the assembled capture machine from a framed configuration equals
the frame of one step of the timed interpreter, including the final halting
transition.
\end{lemma}

\begin{lemma}[Framed runs]
\lean{Turing.timedFrame_run}
\label{Turing.timedFrame_run}
\leanok
\uses{Turing.Cfg, Turing.FinTM, Turing.MultiTapeTM.runFrom,
  Turing.MultiTapeTM.runFrom_succ_eq_step', Turing.TimedControl,
  Turing.timedCaptureTM, Turing.timedFrame, Turing.timedFrame_step,
  Turing.timedInterpreter, Turing.timedUniversalTM}
\difficulty{3}
For every $t$, running the assembled capture machine $t$ steps from a framed
configuration equals the frame of the $t$-step timed-interpreter run, with the
inactive frame unchanged.
\end{lemma}

\begin{lemma}[The captured entry is a frame]
\lean{Turing.timedCaptured_frame}
\label{Turing.timedCaptured_frame}
\leanok
\uses{Turing.Cfg, Turing.FinTM, Turing.FinTM.bufferTape,
  Turing.FinTM.tapeBlocks, Turing.timedCaptureCfg, Turing.timedCapturedCfg,
  Turing.timedCut_InterpreterInitial, Turing.timedFive, Turing.timedFrame,
  Turing.timedLift, Turing.timedSix, Turing.timedUniversalTM,
  Turing.universalFour}
\difficulty{4}
Let a halted configuration of the base machine hold, on the designated clock
tape, the buffer of a clock word with head at the word's length.  Then its
interpreter-entry capture configuration equals the frame --- over the base
configuration's tapes and heads --- of the six-tape lift, with that clock
word, of the stopped interpreter's entry configuration on the captured table
(the base configuration's accumulated output) with the same parked input head.
\end{lemma}

\begin{definition}[The timed universal machine]
\lean{Turing.timedUniversalTM}
\label{Turing.timedUniversalTM}
\leanok
\uses{Turing.EffectiveMachineCode, Turing.FinTM, Turing.timedCanonClock,
  Turing.timedCanonTM, Turing.timedCaptureTM}
The complete timed universal machine of an effective coding scheme: the
capture wrapper of the composed clock-preserving canonizer, with its retained
clock lane.  It has finite control and finitely many tapes.
\end{definition}

\begin{definition}[Code-only startup bound]
\lean{Turing.timedStartupBound}
\label{Turing.timedStartupBound}
\leanok
\uses{Turing.CodeTM.serialize, Turing.EffectiveMachineCode}
The part of the timed universal machine's startup cost that depends only on
the code string $\alpha$: three times $|\alpha|$, plus the canonizer's time
bound at $|\alpha|$, plus the length of the decoded machine's serialization,
plus twice the width of its binary state count, plus twice its initial state
index, plus $16$.
\end{definition}

\begin{lemma}[The header lies inside the serialization]
\lean{Turing.timed_header_bound}
\label{Turing.timed_header_bound}
\leanok
\uses{Turing.CodeTM, Turing.CodeTM.serialize, Turing.universal_pair_length,
  Turing.universal_serialization_header}
\difficulty{2}
For every coded machine, the header endpoint --- twice the width of the binary
state count, plus $2$, plus the initial state index, plus $1$ --- is at most
the length of the machine's serialization.
\end{lemma}

\begin{lemma}[Exact startup of the timed machine]
\lean{Turing.timed_initialized}
\label{Turing.timed_initialized}
\leanok
\uses{Turing.CodeTM.serialize, Turing.EffectiveMachineCode,
  Turing.MultiTapeTM.initCfg, Turing.MultiTapeTM.runFrom,
  Turing.MultiTapeTM.runFrom_add, Turing.pairEncode, Turing.timedCanonClock,
  Turing.timedCanonTM, Turing.timedCanon_complete, Turing.timedCaptureTM,
  Turing.timedCapture_start, Turing.timedCaptured_frame,
  Turing.timedCut_InterpreterBase, Turing.timedCut_InterpreterInitial,
  Turing.timedCut_Interpreter_initialize, Turing.timedFrame,
  Turing.timedFrame_run, Turing.timedLift, Turing.timedStartupBound,
  Turing.timedUniversalTM, Turing.timed_replay, Turing.universalEvalCfg,
  Turing.universalFour, Turing.universalInputPos,
  Turing.universalSimulationCfg, Turing.universal_pair_length,
  Turing.universal_serialization_header}
\difficulty{5}
On the nested pair encoding $\langle \langle b, \alpha\rangle, x\rangle$, the
timed universal machine reaches, within $4|b|$ plus the code-only startup
bound steps, a framed configuration over some inactive tapes and heads: the
six-tape lift, with clock word $b$, of the interpreter checkpoint of the
decoded machine's initial configuration on $x$, with the table cursor just
past the serialization header (at twice the binary state-count width, plus
$2$, plus the initial state index, plus $1$).
\end{lemma}

\begin{lemma}[Quadratic cost absorption]
\lean{Turing.timed_cost_bound}
\label{Turing.timed_cost_bound}
\leanok
\difficulty{4}
For all natural numbers $S$, $B$, $t$, $w$, $s$, $d$: if $w \le t$,
$s \le 4w + S$, and $d \le (B + 2w + 8)(t+1)$, then
$s + d \le (S + B + 14)(t+1)^2$.
\end{lemma}

\begin{lemma}[The timed machine computes the bounded answer]
\lean{Turing.timed_computes}
\label{Turing.timed_computes}
\leanok
\uses{Turing.Cfg.init, Turing.EffectiveMachineCode,
  Turing.FinTM.ComputesInTime, Turing.FinTM.ComputesInTime.mono,
  Turing.FinTM.computesInTime_iff, Turing.MultiTapeTM.initCfg,
  Turing.MultiTapeTM.runFrom, Turing.MultiTapeTM.runFrom_add,
  Turing.pairEncode, Turing.timedAnswer, Turing.timedCanonClock,
  Turing.timedCanonTM, Turing.timedCaptureTM, Turing.timedFrame,
  Turing.timedFrame_run, Turing.timedStartupBound, Turing.timedUniversalTM,
  Turing.timedValue_bits, Turing.timed_bits_length, Turing.timed_cost_bound,
  Turing.timed_header_bound, Turing.timed_initialized,
  Turing.timed_interpret_finishes, Turing.universalBlockBound}
\difficulty{4}
For every effective coding scheme, code string $\alpha$, input $x$, and
deadline $t$, the timed universal machine on the nested pair encoding of the
binary representation of $t$, the code $\alpha$, and the input $x$ halts with
output the deadline-inclusive answer of the decoded machine's initial
configuration at budget $t$, within
$(\text{startup bound} + \text{block bound} + 14)\,(t+1)^2$ steps, where both
bounds depend only on $\alpha$.
\end{lemma}

\begin{theorem}[The time-bounded universal machine]
\lean{Turing.timed_universal}
\label{Turing.timed_universal}
\leanok
\uses{Turing.CodeTM.toFinTM, Turing.EffectiveMachineCode, Turing.FinTM,
  Turing.FinTM.ComputesInTime, Turing.FinTM.computesInTime_iff,
  Turing.MultiTapeTM.initCfg, Turing.MultiTapeTM.runFrom, Turing.pairEncode,
  Turing.timedAnswer, Turing.timedStartupBound, Turing.timedUniversalTM,
  Turing.timed_computes, Turing.universalBlockBound}
\difficulty{4}
For every effective machine-coding scheme there is a single finite-tape
machine $U$ over $\{0,1\}$ such that for every code string $\alpha$ there is a
constant $C$ with the following two properties for every input $x$ and every
deadline $t \in \mathbb{N}$.  If the machine denoted by $\alpha$ halts on $x$
within $t$ steps with output $w$, then $U$ on the nested pair encoding of the
binary representation of $t$, the code $\alpha$, and the input $x$ halts with
output $1$ followed by $w$ within $C (t+1)^2$ steps.  If instead the denoted
machine produces no output on $x$ within $t$ steps, then $U$ halts with output
the single timeout tag $0$ within the same bound.  Halting is checked after
every simulated transition including the $t$-th, so halting exactly at the
deadline counts as success.
\end{theorem}
